\subsection{The economic equilibrium}\label{\refsec:EE}

\smalltitle{With log preferences.}
Given the paths of taxes $\tau_t$ and pension characteristics $\theta_t, \epsilon_t$ and an initial level of savings $s_0$, the economic equilibrium needs no numerical root finding, since $B_{t+1}^i = \beta_{t,i}$:
\begin{itemize}
    \item compute $\Gamma_{s,t}$ from \eqref{\refeq:auxiliary:Gammas} for $t<T$;
    \item compute $\Theta_{h,t}$ from \eqref{\refeq:auxiliary:Thetah} for $t<T$ and \eqref{\refeq:auxiliary:ThetahT} for $t=T$;
    \item compute $\Theta_{s,t}$ from \eqref{\refeq:auxiliary:Thetas} for $t<T$;
    \item iterate \eqref{\refeq:auxiliary:s} forward from $s_0$ for $t<T$, with $s_T = 0$.
\end{itemize}
These core variables give individual labour supply and savings, factor prices, pension benefits and consumption as in \cref{\refmodelsec:EE}.


\smalltitle{With CRRA preferences.}
With the same inputs, the core of the economic equilibrium is the paths of $\Gamma_{s,t}$, $h_t$ and $s_t$ for $t<T$, identified by a square system of nonlinear equations:
\begin{itemize}
    \item labour supply, \eqref{\refeq:auxiliary:h}, with $\Theta_{h,t}$ from \eqref{\refeq:auxiliary:Thetah} for $t<T$ and \eqref{\refeq:auxiliary:ThetahT} for $t=T$;
    \item savings, \eqref{\refeq:auxiliary:sFromH}, for $t<T$, with $s_T = 0$;
    \item $\Gamma_{s,t}$ from \eqref{\refeq:auxiliary:Gammas} for $t<T$, with $B_{t+1}^i(s_t, h_{t+1})$ from \eqref{\refeq:auxiliary:B}.
\end{itemize}
A standard root finder solves the system, started from the previous solution or, without one, from the log solution, which is exact at $\rho = 1$. The rest of the equilibrium follows as in the log case.


\smalltitle{Initializing with steady state savings.}
By default the initial savings $s_0$ are the steady state at the policy and the parameters of the first period ($t=1$). With log preferences $B^i = \beta_i$, so $\Gamma_s$ is the closed form \eqref{\refeq:auxiliary:Gammas}, and the steady state follows:
\begin{subequations}\label{\refeq:steadystate_LOG}
    \begin{align}
        \Gamma_s =& \frac{1}{1+\xi}\frac{\sum_i \frac{\gamma_{i}\eta_{i}^{1+\xi}}{X_{i}^{\xi}}\frac{B^i}{1+B^i}}{1+\frac{1-\alpha}{\alpha}\frac{p\tau}{\kappa}\left(\theta+(1-\theta)\sum_i\frac{\gamma_{i}}{1+B^i}\right)} \label{\refeq:steadystate_LOG:Gammas}\\
         s^* =& \Gamma_h^{1+\xi}\left[\left(\frac{(1-\alpha)(1-\tau)}{\Gamma_h-\frac{1-\alpha}{\alpha}\frac{p\theta\tau}{\kappa}\Gamma_s}\right)^{1+\xi}\frac{\Gamma_s^{1+\alpha\xi}}{\nu^{\alpha(1+\xi)}}\right]^{\frac{1}{1-\alpha}} \label{\refeq:steadystate_LOG:s}
    \end{align}
\end{subequations}

With CRRA preferences $s^*$ is found numerically, as the $\Gamma_s$ that solves the scalar root problem
\begin{subequations}\label{\refeq:steadystate_CRRA}
\begin{align}
    &\Gamma_s(1+\xi)\left[1+\frac{1-\alpha}{\alpha}\frac{p\tau}{\kappa}\left(\theta+(1-\theta)\sum_i\frac{\gamma_i}{1+B^i(\Gamma_s)}\right)\right] = \sum_i\frac{\gamma_i\eta_i^{1+\xi}}{X_i^{\xi}} \frac{B^i(\Gamma_s)}{1+B^i(\Gamma_s)}, \label{\refeq:steadystate_CRRA:Gammas} \\
    &\text{where }B^i(\Gamma_s) = \left(\beta_i\right)^{\rho} \left(\frac{\alpha}{p}\frac{\Gamma_h-\frac{1-\alpha}{\alpha}\frac{p\theta\tau}{\kappa}\Gamma_s}{(1-\alpha)(1-\tau)}\frac{\nu}{\Gamma_s}   \right)^{\rho-1}. \label{\refeq:steadystate_CRRA:Bi}
\end{align}
\end{subequations}
The root is bracketed on both sides. Since $B^i/(1+B^i)<1$ and the bracket on the left of \eqref{\refeq:steadystate_CRRA:Gammas} exceeds one, $\Gamma_s<\Gamma_h/(1+\xi)$, which is $0.77$ at $\Gamma_h = 1$ and $\xi = 0.3$. And $\Theta_{h}$ and $B^i(\Gamma_s)$ are defined only below $\bar{\Gamma}_s\equiv\alpha\kappa\Gamma_h/\left((1-\alpha)p\theta\tau\right)$, where the denominator of \eqref{\refeq:auxiliary:Thetah} vanishes; the cap is infinite at $\theta\tau = 0$ and scales with $\alpha/(1-\alpha)$. We therefore apply Brent's method on $\left(10^{-6}, \min(0.75, 0.99\,\bar{\Gamma}_s)\right)$, and double the upper end, up to $0.99\,\bar{\Gamma}_s$, where that interval does not bracket a sign change.
